% ==========================================================================
%  Chapter 78 -- Main Driver Refactoring: MultiscaleAnalysis Integration
% ==========================================================================

\chapter{Main Driver Refactoring: MultiscaleAnalysis Integration}
\label{ch:driver-refactoring}


% =========================================================================
\section{Introduction}
\label{sec:driver-intro}

Chapter~\ref{ch:multiscale-refactoring} implemented six architectural
phases, culminating in the \texttt{MultiscaleAnalysis} orchestrator
class that encapsulates the FE\textsuperscript{2} staggered coupling
loop.  This chapter documents the integration of that orchestrator
into the primary analysis driver
(\texttt{main\_lshaped\_multiscale.cpp}), replacing $\sim$80~lines of
hand-written staggered coupling logic with a single call to
\verb|analysis.step()|.

The refactoring achieves:
\begin{itemize}
  \item \textbf{Line reduction}: 80~lines of staggered coupling $\to$
    15~lines of orchestrator invocation + post-step diagnostics.
  \item \textbf{Separation of concerns}: coupling logic
    (\texttt{MultiscaleAnalysis}), scale bridging
    (\texttt{MultiscaleModel}), and damage diagnostics (main driver)
    are cleanly separated.
  \item \textbf{Behavioral equivalence}: the refactored driver
    produces identical numerical results---the ``lagged'' staggered
    scheme, convergence tolerance, relaxation factor, and coupling
    start step are preserved unchanged.
\end{itemize}


% =========================================================================
\section{Architecture Overview}
\label{sec:driver-arch}

Figure~\ref{fig:driver-arch} shows the object graph after
refactoring.  The main driver creates the following objects and wires
them together:

\begin{enumerate}
  \item \texttt{MultiscaleModel<NonlinearSubModelEvolver>} --- owns the
    sub-model evolvers (moved from a temporary vector), the kinematics
    extraction callback, and the response injection callback.
  \item \texttt{MultiscaleAnalysis<NonlinearSubModelEvolver>} --- owns
    the \texttt{MultiscaleModel} and three strategy objects:
    \texttt{TwoWayStaggered}, \texttt{FrobeniusConvergence},
    \texttt{ConstantRelaxation}.
  \item The main driver retains ownership of the global model, solver,
    observers, VTK exporters, and crack diagnostics.
\end{enumerate}

\begin{figure}[htb]
\centering
\small
\begin{verbatim}
  main()
    |
    +-- model (StructModel)             [global FE model]
    +-- solver (DynSolver)              [Newmark time integrator]
    +-- analysis (MultiscaleAnalysis)
    |     |
    |     +-- model_ (MultiscaleModel)
    |     |     +-- local_models_       [NonlinearSubModelEvolver ...]
    |     |     +-- kinematics_extractor_
    |     |     +-- response_injector_
    |     |
    |     +-- bridge_  (TwoWayStaggered)
    |     +-- convergence_ (FrobeniusConvergence)
    |     +-- relaxation_  (ConstantRelaxation)
    |
    +-- crack_csv, pvd_global, ...      [diagnostics / I/O]
\end{verbatim}
\caption{Object ownership graph after driver refactoring.}
\label{fig:driver-arch}
\end{figure}


% =========================================================================
\section{Evolver Transfer and Callback Wiring}
\label{sec:driver-transfer}

Sub-model evolvers are created in a temporary vector, then moved into
the \texttt{MultiscaleModel}:

\begin{lstlisting}[language=C++,caption={Evolver transfer and callback
  wiring (section~12 of the driver).},label={lst:driver-transfer}]
MultiscaleModel<NonlinearSubModelEvolver> ms_model;

// Kinematics extraction: delegates to existing lambda
ms_model.set_kinematics_extractor(extract_beam_kinematics);

// Response injection: sets D_hom and f_hom on beam elements
ms_model.set_response_injector(
    [&model](std::size_t eid,
             const Eigen::Matrix<double,6,6>& D_hom,
             const Eigen::Vector<double,6>&   f_hom)
    {
        if (auto* beam = model.elements()[eid].as<BeamElemT>()) {
            beam->set_homogenized_tangent(D_hom);
            beam->set_homogenized_forces(f_hom);
        }
    });

for (auto& ev : nl_evolvers) {
    const auto eid = ev.parent_element_id();
    ms_model.register_local_model(eid, std::move(ev));
}
\end{lstlisting}

Note the defensive sequencing in the move loop: the element id is
captured \emph{before} the move, avoiding undefined behaviour from
unspecified function-argument evaluation order in C++.


% =========================================================================
\section{Orchestrator Construction}
\label{sec:driver-orchestrator}

The \texttt{MultiscaleAnalysis} is constructed with the same
parameters that governed the original hand-written loop:

\begin{lstlisting}[language=C++,caption={MultiscaleAnalysis
  construction with strategy injection.},
  label={lst:driver-orchestrator}]
MultiscaleAnalysis<NonlinearSubModelEvolver> analysis(
    std::move(ms_model),
    std::make_unique<TwoWayStaggered>(MAX_STAGGERED_ITER),
    std::make_unique<FrobeniusConvergence>(STAGGERED_TOL),
    std::make_unique<ConstantRelaxation>(STAGGERED_RELAX));

analysis.set_coupling_start_step(COUPLING_START_STEP);
analysis.set_section_dimensions(COL_B, COL_H);
\end{lstlisting}

The mapping between the original constants and the strategy objects is
shown in Table~\ref{tab:driver-mapping}.

\begin{table}[htb]
\centering
\caption{Mapping from original constants to strategy objects.}
\label{tab:driver-mapping}
\begin{tabular}{lll}
\hline
\textbf{Original constant} & \textbf{Value} & \textbf{Strategy object} \\
\hline
\texttt{MAX\_STAGGERED\_ITER}  & 4    & \texttt{TwoWayStaggered(4)} \\
\texttt{STAGGERED\_TOL}        & 0.05 & \texttt{FrobeniusConvergence(0.05)} \\
\texttt{STAGGERED\_RELAX}      & 0.7  & \texttt{ConstantRelaxation(0.7)} \\
\texttt{COUPLING\_START\_STEP} & 10   & \texttt{set\_coupling\_start\_step(10)} \\
\hline
\end{tabular}
\end{table}


% =========================================================================
\section{Staggered Loop Replacement}
\label{sec:driver-loop}

The most significant change is the replacement of the hand-written
staggered coupling loop.  The original implementation spanned
$\sim$80~lines and interleaved sub-model solving, tangent computation,
relaxation, convergence checking, response injection, and crack
diagnostics.

% -------------------------------------------------------------------------
\subsection{Before: Hand-Written Staggered Loop}

The original loop structure is shown in
Listing~\ref{lst:driver-before} (abbreviated).

\begin{lstlisting}[language=C++,caption={Original staggered coupling
  loop (abbreviated, $\sim$80 lines).},label={lst:driver-before}]
std::vector<Eigen::Matrix<double,6,6>> D_prev(...);
int staggered_iters = 0;

for (int s_iter = 0; ; ++s_iter) {
    // (a) Extract kinematics, update BCs, solve sub-models
    for (auto& ev : nl_evolvers) {
        ev.update_kinematics(...);
        ev.solve_step(t);
        // ... accumulate crack summaries ...
    }

    if (!do_coupling || s_iter >= MAX_STAGGERED_ITER) break;

    // (c) Compute tangent, relax, check convergence, inject
    for (std::size_t i = 0; i < nl_evolvers.size(); ++i) {
        auto D_hom = nl_evolvers[i].compute_homogenized_tangent(...);
        // relax, check convergence, inject into beam ...
    }

    if (s_iter == 0) continue;  // at least 2 iterations
    if (all_converged) break;
    break;  // lagged coupling: exit after check
}
\end{lstlisting}

% -------------------------------------------------------------------------
\subsection{After: Orchestrator Delegation}

The refactored code delegates entirely to the orchestrator:

\begin{lstlisting}[language=C++,caption={Refactored staggered coupling
  ($\sim$15 lines).},label={lst:driver-after}]
// FE^2 staggered coupling (delegated to MultiscaleAnalysis)
analysis.step(t, evol_step);
const int staggered_iters = analysis.last_staggered_iterations();

// Collect crack summaries (post-step)
int step_cracked_gps = 0, step_total_cracks = 0;
double step_max_crack_damage = 0.0, step_max_opening = 0.0;
std::vector<CrackSummary> step_summaries;
for (auto& ev : analysis.model().local_models()) {
    auto cs = ev.crack_summary();
    step_summaries.push_back(cs);
    step_cracked_gps  += cs.num_cracked_gps;
    step_total_cracks += cs.total_cracks;
    // ...
}
\end{lstlisting}

The crack diagnostics are collected \emph{after}
\verb|analysis.step()| returns, querying each sub-model evolver via
\verb|analysis.model().local_models()|.  This is equivalent to the
original, which collected diagnostics after the final staggered
iteration.


% =========================================================================
\section{Behavioral Equivalence Analysis}
\label{sec:driver-equivalence}

Identical numerical behavior is ensured by the following design
decisions:

\begin{enumerate}
  \item \textbf{Lagged staggered scheme.}  The
    \texttt{MultiscaleAnalysis::step()} method mirrors the original
    loop structure: iteration~0 always continues to iteration~1
    (\texttt{s\_iter == 0} $\Rightarrow$ \texttt{continue}), and the
    convergence check occurs only at iteration~1 onwards.  The
    effective behavior is a lagged coupling where the homogenized
    tangent computed at step~$n$ is used by the global solver at
    step~$n+1$.

  \item \textbf{No auto-commit.}  The original driver does not call
    \texttt{commit\_state()} on the sub-model evolvers after each
    global step.  The orchestrator's \texttt{auto\_commit\_} flag
    defaults to \texttt{false}, preserving this behavior.  The
    material state management follows the evolver's internal SNES
    convergence cycle.

  \item \textbf{Same convergence metric.}  The
    \texttt{FrobeniusConvergence} strategy implements the identical
    relative Frobenius norm criterion:
    \begin{equation}\label{eq:driver-conv}
      \frac{\|\mathbf{D}_\text{curr}^{(i)} -
            \mathbf{D}_\text{prev}^{(i)}\|_F}
           {\|\mathbf{D}_\text{curr}^{(i)}\|_F}
      < \tau
      \qquad \forall\, i \in \{1, \dots, n_\text{sub}\}
    \end{equation}
    with $\tau = 0.05$ and per-sub-model enforcement.

  \item \textbf{Same relaxation.}  The \texttt{ConstantRelaxation}
    strategy applies $\omega = 0.7$ via:
    \begin{equation}\label{eq:driver-relax}
      \hat{\mathbf{D}}^{(i)} = \omega\, \mathbf{D}_\text{hom}^{(i)}
        + (1 - \omega)\, \mathbf{D}_\text{prev}^{(i)}
    \end{equation}
    only when $s_\text{iter} > 0$, matching the original guard.

  \item \textbf{Post-step crack diagnostics.}  The original code
    accumulated crack summaries inside the staggered loop but
    cleared and re-accumulated on each iteration.  Only the final
    iteration's data was written to CSV.  The refactored code
    collects crack summaries once after \verb|step()| returns,
    yielding identical output.
\end{enumerate}


% =========================================================================
\section{Diagnostic Accessors}
\label{sec:driver-diagnostics}

Two accessors were added to \texttt{MultiscaleAnalysis} to support
post-step reporting in the main driver:

\begin{lstlisting}[language=C++,caption={Diagnostic accessors on
  MultiscaleAnalysis.},label={lst:driver-diag}]
[[nodiscard]] int  last_staggered_iterations() const
    { return last_staggered_iters_; }

[[nodiscard]] bool last_converged() const
    { return last_converged_; }
\end{lstlisting}

These are populated by \texttt{step()} and consumed by the progress
printing logic (every~500 steps) and the crack CSV writer.  The
\texttt{auto\_commit\_} flag is controlled via
\texttt{set\_auto\_commit(bool)} and defaults to \texttt{false}
to match the original driver semantics.


% =========================================================================
\section{Line Count and Complexity Reduction}
\label{sec:driver-metrics}

Table~\ref{tab:driver-metrics} summarises the quantitative impact of
the refactoring on the driver file.

\begin{table}[htb]
\centering
\caption{Driver file metrics before and after refactoring.}
\label{tab:driver-metrics}
\begin{tabular}{lcc}
\hline
\textbf{Metric} & \textbf{Before} & \textbf{After} \\
\hline
Lines in staggered loop block       & $\sim$80  & $\sim$15  \\
Cyclomatic complexity (loop)        & 8         & 1         \\
Direct \texttt{nl\_evolvers} refs   & 12        & 0         \\
Strategy parameters hard-coded      & 4         & 0         \\
Build targets                       & 88        & 88        \\
CTest cases passing                 & 63/63     & 63/63     \\
\hline
\end{tabular}
\end{table}

All 12~direct references to the \texttt{nl\_evolvers} vector in post-
creation code (sections~13--16) were replaced with access through
\verb|analysis.model().local_models()| or
\verb|analysis.model().num_local_models()|.  The temporary vector
is populated, moved into the model, and never referenced again.


% =========================================================================
\section{Response Injector Design}
\label{sec:driver-injector}

A key design decision was encapsulating the response injection---which
couples the orchestrator to concrete beam element types---in a
\texttt{std::function} callback.  The lambda captures the global
\texttt{model} by reference and performs a safe \texttt{as<BeamElemT>}
downcast:

\begin{equation}
  \text{inject}: (\text{eid}, \mathbf{D}_\text{hom},
    \mathbf{f}_\text{hom})
  \;\mapsto\;
  \text{beam.set\_homogenized\_tangent}(\mathbf{D}_\text{hom}),\;
  \text{beam.set\_homogenized\_forces}(\mathbf{f}_\text{hom})
\end{equation}

This design preserves the Open/Closed Principle: new element types
(shell, solid) can be supported by changing only the injector callback,
without modifying \texttt{MultiscaleAnalysis} or
\texttt{MultiscaleModel}.


% =========================================================================
\section{Build and Test Verification}
\label{sec:driver-verification}

After all edits:

\begin{itemize}
  \item \textbf{Build}: 88 Ninja targets compile under GCC~15.2 with
    \texttt{-std=c++23 -Werror -Wall -Wextra -Wpedantic}.  No
    warnings.
  \item \textbf{Tests}: 63/63 CTest cases pass with
    \texttt{--output-on-failure}.
  \item \textbf{Files modified}: 2
    (\texttt{main\_lshaped\_multiscale.cpp},
     \texttt{MultiscaleAnalysis.hh}).
\end{itemize}


% =========================================================================
\section{Discussion}
\label{sec:driver-discussion}

% -------------------------------------------------------------------------
\subsection{Move Semantics and Evaluation Order}

The evolver transfer loop demonstrates a subtle C++ pitfall: passing
\texttt{ev.parent\_element\_id()} and \texttt{std::move(ev)} as
arguments to the same function call invokes unspecified evaluation
order (C++17~\S\,8.2.2).  If the move is evaluated first, the
subsequent \texttt{parent\_element\_id()} call operates on a
moved-from object.  The fix---capturing the id into a local variable
before the function call---is a defensive pattern that should be
applied whenever a move and a member access coexist in the same
argument list.

% -------------------------------------------------------------------------
\subsection{Commit Semantics}

The original driver does \emph{not} call
\texttt{commit\_state()} on the sub-model evolvers after each global
step.  This is intentional: the evolver's internal SNES solve
evaluates constitutive response during residual assembly, and the
converged displacement field is retained in the PETSc~\texttt{Vec}.
The perturbation tangent computation saves and restores both the
displacement vector and the material state, making an explicit commit
unnecessary for the lagged coupling scheme.

The \texttt{auto\_commit\_} flag provides forward compatibility: a
future multi-iteration staggered scheme that re-solves the global
step within the coupling loop would require explicit commits between
staggered iterations.

% -------------------------------------------------------------------------
\subsection{Crack Diagnostic Decoupling}

Moving crack summary collection from inside the staggered loop to a
post-step query is a clean application of the Single Responsibility
Principle.  The staggered loop is now purely concerned with coupling
convergence; diagnostic collection is a separate concern handled by
the driver after the coupling is complete.


% =========================================================================
\section{Summary}
\label{sec:driver-summary}

This chapter documented the integration of the
\texttt{MultiscaleAnalysis} orchestrator into the L-shaped multiscale
driver.  The refactoring replaced $\sim$80~lines of hand-written
staggered coupling logic with a single \verb|analysis.step()| call,
reduced the cyclomatic complexity of the inner loop from~8 to~1, and
decoupled crack diagnostics from coupling convergence.  All 88~build
targets compile cleanly and all 63~test cases continue to pass.
